% from ega4-2-fr.tex
\oldpage[IV-2]{205}

\begin{lemma}[7.5.1.1]
\label{IV.2.s7.5.1.1-fr}
Soit $\mathcal C$ une sous-catégorie pleine de la catégorie des anneaux
locaux noethériens, telle que tout anneau quotient d'un anneau de
$\mathcal C$ appartienne encore à $\mathcal C$. Soit $\boldsymbol R(C)$ une
propriété telle que si $C\in\mathcal C$ et si $t$ est un élément régulier de
l'idéal maximal de $C$ tel que $\boldsymbol R(C/tC)$ soit vraie, alors
$\boldsymbol R(C)$ est vraie.

Cela étant, soient $C$ un anneau local régulier, $k$ son corps résiduel,
$D$ un anneau local appartenant à $\mathcal C$,
$\varphi:C\to D$ un homomorphisme local faisant de $D$ un $C$-module plat.
Alors, si $\boldsymbol R(D\otimes_C k)$ est vraie, il en est de même de
$\boldsymbol R(D)$.
\end{lemma}

\begin{proof}
Raisonnons par récurrence sur $n=\dim(C)$, le lemme étant vrai par hypothèse
si $n=0$, puisqu'alors $C=k$. Soit $t$ un élément de l'idéal maximal
$\mathfrak m$ de $C$ n'appartenant pas à $\mathfrak m^2$; on sait alors que
$C/tC$ est régulier et que $\dim(C/tC)=n-1$ ((0, 17.1.8) et (0, 16.3.4)).
Comme
\[
(D/tD)\otimes_{C/tC}k=D\otimes_C k,
\]
et que $D/tD\in\mathcal C$, l'hypothèse de récurrence montre que
$\boldsymbol R(D/tD)$ est vraie; en outre, $t$ est $C$-régulier, donc aussi
$D$-régulier par platitude (0$_{\mathrm I}$, 6.3.4); l'hypothèse faite sur
$\boldsymbol R$ montre donc que $\boldsymbol R(D)$ est vraie.
\end{proof}

Pour terminer la démonstration de (7.5.1), il suffit ici d'appliquer le
lemme (7.5.1.1) en prenant pour $\mathcal C$ la catégorie des anneaux locaux
des spectres d'anneaux locaux complets, tenant compte de l'hypothèse
$(\mathrm R_{\mathrm{IV}})$, et prenant pour $C$ l'anneau
$A_{\mathfrak p}$, pour $D$ l'anneau $B_{\mathfrak q}$.
\end{proof}

\begin{corollary}[7.5.2]
\label{IV.2.s7.5.2-fr}
Soient $A$ un anneau local noethérien, $\mathfrak m$ l'idéal maximal de $A$,
$k=A/\mathfrak m$ son corps résiduel, $B$ un anneau local noethérien,
$\varphi:A\to B$ un homomorphisme local tel que :
\begin{enumerate}
\item[\emph{(i)}] Le corps résiduel de $B$ est une extension finie de $k$.
\item[\emph{(ii)}] $B$ est un $A$-module plat.
\end{enumerate}
Soit d'autre part $\boldsymbol R(C)$ une propriété vérifiant les conditions
$(\mathrm R_{\mathrm{II}})$ (7.3.10), $(\mathrm R_{\mathrm{III}})$
(7.3.18) et $(\mathrm R_{\mathrm{IV}})$ (7.5.1). Enfin, soit
$\boldsymbol R'(C)$ une propriété vérifiant la condition suivante :
\begin{quote}
$(\mathrm R_{\mathrm V})$ Si $C$, $D$ sont deux anneaux locaux noethériens,
$\psi:C\to D$ un homomorphisme local tel que
\[
g={}^a\psi:\operatorname{Spec}(D)\to\operatorname{Spec}(C)
\]
soit un $\boldsymbol P$-morphisme, alors la propriété
$\boldsymbol R'(C)$ implique $\boldsymbol R(D)$.
\end{quote}
Supposons que le morphisme canonique
$\operatorname{Spec}(\widehat A)\to\operatorname{Spec}(A)$ soit un
$\boldsymbol P'$-morphisme, où $\boldsymbol P'$ se définit à partir de
$\boldsymbol R'$ comme $\boldsymbol P$ à partir de $\boldsymbol R$ dans
(7.5.0).

Cela étant, soit
$f={}^a\varphi:\operatorname{Spec}(B)\to\operatorname{Spec}(A)$, et
supposons que la propriété
\[
\boldsymbol P(\operatorname{Spec}(\widehat B
\otimes_{\widehat A}k),k)
\]
soit vraie. Alors, pour tout $x\in\operatorname{Spec}(A)$, la propriété
$\boldsymbol P(f^{-1}(x),k(x))$ est vraie (autrement dit, $f$ est un
$\boldsymbol P$-morphisme).
\end{corollary}

\begin{proof}
Procédons encore en deux étapes :

\emph{I) Réduction à l'étude des anneaux locaux de la fibre générique.} ---
Appliquons encore le lemme (7.3.16.2); la seule différence avec le
raisonnement de I) dans (7.5.1) est qu'ici l'anneau $A'$ est seulement un
anneau semi-local intègre, mais non nécessairement local; tout idéal maximal
$\mathfrak n'$ de $B'$ est alors au-dessus d'un idéal maximal $\mathfrak m'$
de $A'$; comme tout anneau local de
$\operatorname{Spec}(B'\otimes_{A'}K')$ est un anneau local de
$\operatorname{Spec}(B')$, donc de l'un des
$\operatorname{Spec}(B'_{\mathfrak n'})$ (au-dessus du point générique de
$\operatorname{Spec}(A'_{\mathfrak m'}))$ on est ramené à prouver que les
anneaux locaux de
\[
\operatorname{Spec}(B'_{\mathfrak n'}
\otimes_{A'_{\mathfrak m'}}K')
\]
possèdent la propriété $\boldsymbol R$. Or, la définition de
$\boldsymbol P'$ implique que si $\boldsymbol P'(Z,k)$ est vraie, il en est
de même de $\boldsymbol P'(Z\otimes_k k',k')$ pour toute extension finie
$k'$ de $k$; le même raisonnement que dans (7.3.7) prouve que les fibres
formelles de $A'_{\mathfrak m'}$ possèdent la propriété $\boldsymbol P'$.
On voit comme dans la partie I) de la démonstration de (7.5.1) que
$A'_{\mathfrak m'}$ et $B'_{\mathfrak n'}$ vérifient les conditions (i) et
(ii) de (7.5.2); d'autre part, si $k'$ est le corps résiduel de
$A'_{\mathfrak m'}$, la propriété
\[
\boldsymbol P(\operatorname{Spec}(\widehat{B'_{\mathfrak n'}}
\otimes_{\widehat{A'_{\mathfrak m'}}}k'),k')
\]
est vraie : en effet, le complété de $A'_{\mathfrak m'}$ est l'un des
anneaux locaux composants de
$\widehat{A'}=\widehat A\otimes_A A'$ et le complété de
$B'_{\mathfrak n'}$ l'un des anneaux locaux composants de
\[
\widehat{B'}=\widehat B\otimes_B B'
=\widehat B\otimes_A A'
=\widehat B\otimes_{\widehat A}\widehat{A'};
\]
le raisonnement de I), dans (7.5.1), prouve donc notre assertion.

On peut donc se borner au cas où $A$ est en outre supposé intègre, et à
prouver que, si $K$ est le corps des fractions de $A$, les anneaux locaux de
$\operatorname{Spec}(B\otimes_A K)$ possèdent la propriété
$\boldsymbol R$.

\emph{II) Réduction au cas où $A$ et $B$ sont complets.} --- Soient $\xi$
le point générique de $\operatorname{Spec}(A)$, $y$ un point quelconque de
$f^{-1}(\xi)$; il s'agit de prouver que
$\boldsymbol R(\mathscr O_y)$ est vraie; compte tenu de la condition
$(\mathrm R_{\mathrm{II}})$, il suffira de voir qu'il existe dans
$\operatorname{Spec}(\widehat B)$ un point $z$ au-dessus de $y$ et tel que
$\boldsymbol R(\mathscr O_z)$ soit vraie. Or, il résulte de l'hypothèse
faite sur $\operatorname{Spec}(\widehat B\otimes_{\widehat A}k)$ et de
(7.5.1) que le morphisme
\enlargethispage{5\baselineskip}
\[
{}^a\widehat\varphi:\operatorname{Spec}(\widehat B)
\to\operatorname{Spec}(\widehat A)
\]
est un $\boldsymbol P$-morphisme. D'autre part, pour tout point
$z\in\operatorname{Spec}(\widehat B)$ au-dessus de $y$ (il en existe,
puisque $\widehat B$ est un $B$-module fidèlement plat), l'image $x$ de $z$
dans $\operatorname{Spec}(\widehat A)$ appartient à la fibre
$h^{-1}(\xi)$ pour le morphisme canonique
$h:\operatorname{Spec}(\widehat A)\to\operatorname{Spec}(A)$; en vertu de
l'hypothèse faite sur $A$, la propriété
$\boldsymbol R'(\mathscr O_x)$ est vraie; donc en vertu de
$(\mathrm R_{\mathrm V})$, $\boldsymbol R(\mathscr O_z)$ est vraie.
C.Q.F.D.
\end{proof}

\begin{examples}[7.5.3]
\label{IV.2.s7.5.3-fr}
Considérons les propriétés $\boldsymbol R(C)$ suivantes ($C$ étant un anneau
local noethérien) :
\begin{enumerate}
\item[\emph{(i)}] (aussi notée $(\mathrm i_n)$) $\operatorname{coprof}(C)\leq n$.
